\documentclass[10pt,a4paper]{amsart}
\usepackage[margin=19mm]{geometry}
\usepackage{amsmath,amssymb,mathtools}
\usepackage[scr=rsfs,cal=euler]{mathalfa}
\usepackage{xfrac}
\usepackage[unicode,hidelinks,bookmarks=false]{hyperref}
\newcommand{\ie}{i.e.\ }
\newcommand{\cf}{cf.\ }
\newcommand{\resp}{respectively\ }
\newcommand{\Subsection}{\subsection}
\providecommand{\nobreakdash}{\nobreak\hbox{-}}
\setlength{\emergencystretch}{2em}
\multlinegap0pt
\usepackage{mathtools}
\usepackage{enumitem}
\usepackage{graphicx}
\newcommand{\uu}{\boldsymbol{\mathfrak{u}}}             % was \boldsymbol{\mathfrak{u}}
\newcommand{\m}{\mathfrak{m}}
\newcommand{\nn}{\mathfrak{n}}
\newcommand{\I}{\mathcal{I}}
\newcommand{\R}{\mathbf{R}}              % was \boldsymbol{\mathrm{R}}
\newcommand{\h}{\mathsf{h}}
\newcommand{\s}{\mathbb{S}}
\newcommand{\Dsf}{\mathsf{D}}
\newcommand{\di}{\mathsf{d}} 
\newcommand{\cs}{\texttt{C}}
\newcommand{\z}{\mathsf{z}}
\newcommand{\Lsf}{\mathsf{L}}
\newcommand{\Msf}{\mathsf{G}}
\newcommand{\F}{\mathbb{F}}
\newcommand{\U}{\mathbf{U}}              % was \boldsymbol{\mathrm{U}}
\newcommand{\A}{\boldsymbol{\mathfrak{U}}}              % was \boldsymbol{\mathrm{a}}
\newcommand{\B}{{\mathfrak{M}}}
\newcommand{\x}{\mathbf{X}}              % was \boldsymbol{\mathrm{X}}
\newcommand{\Q}{\mathbf{Q}}              % was \boldsymbol{\mathrm{Q}}
\newcommand{\Ell}{\mathfrak{L}}
\newcommand{\D}{\mathcal{D}}
\newcommand{\ee}{\mathfrak{e}}

\newcommand{\bu}{\mathbf{U}}             % was \boldsymbol{\mathrm{u}}
\newcommand{\ud}{\,\mathrm{d}}
\newcommand{\e}{{\varepsilon}}
%\renewcommand{\O}{\mathcal{O}}
%\newtheorem{lemma}{Lemma}
%\newtheorem{corollary}{Corollary}

\newcommand{\BB}{\mathbb{B}}            % matrix B
\renewcommand{\AA}{\mathbb{A}}            % matrix A
\newcommand{\II}{\mathbf{I}}            % identity
\newcommand{\xx}{\mathbf{x}}            % state vector
\newcommand{\cc}{\mathbf{c}}            % Fourier coefficient
\newcommand{\ww}{\mathbf{w}}            % left eigenvector
\newcommand{\fp}{\mathfrak{p}}          % shorthand p
\newcommand{\fq}{\mathfrak{q}}          % shorthand q
\newcommand{\fr}{\mathfrak{r}}          % shorthand r
\newcommand{\tr}{\operatorname{tr}}     % trace
\newcommand{\arcosh}{\operatorname{arcosh}} % inverse cosh
\newcommand{\RR}{\mathcal{R}}
\newcommand{\SSS}{\mathcal{S}}

%%%%%%%%%


\newtheorem{theorem}{Theorem}
\newtheorem{lemma}[theorem]{Lemma}
\newtheorem{corollary}[theorem]{Corollary}
\newtheorem{proposition}[theorem]{Proposition}

\newtheorem{remark}[theorem]{Remark}
\title[On the instability of some upward propagating exact nonlinear mountain waves]{On the instability of some upward propagating, exact, nonlinear mountain waves: the unstable set of stratified layers (Theorem 13)}
\author[Christian Puntini]{Christian Puntini}
\date{Source: arXiv:2604.03620v3 (CC BY 4.0)}
\begin{document}
\maketitle
\vspace{1em}
\noindent\textit{Source and licence.} On the instability of some upward propagating, exact, nonlinear mountain waves. Version arXiv:2604.03620v3 (CC BY 4.0), distributed under the Creative Commons Attribution 4.0 International licence, http://creativecommons.org/licenses/by/4.0/, as recorded in the arXiv OAI licence metadata for this exact version and shown on the abstract page https://arxiv.org/abs/2604.03620v3. This item is an editorial re-presentation of the paper's Theorem 13 together with the paper's own proof of it, the retained source-local chain it uses, and the two bibliography entries the retained text cites, transcribed from the paper's own bibliography. A full rights notice is delivered at licenses/arxiv-2604.03620-CC-BY-4.0.txt.
\noindent\textit{Editorial scope.} Counted unit: the paper's real-analyticity and structure theorem for the discriminant $\Delta(\ee,\nu^2)$ of stratified layers, printed as Theorem 13 (source0.tex lines 1040--1050), delivered together with the paper's complete original proof of it (source0.tex lines 1052--1236) as the single primary proof body. No mathematics is written, repaired or re-derived: statement and proof are the paper's own text line for line, in the paper's own notation ($\Delta$, $\ee$, $\nu^2$, $\mathbb M$, $\mathcal S$, $\ee_c$, $\lambda$, $\D$, $\Msf$, $\A$, $\B$). Required context retained uncounted, because the counted proof is the paper's own assembly over it: the isentropic relation (source0.tex lines 255--257); the leading-order WKB amplitude system (lines 369--376); the stratification classification with the linear system for $\boldsymbol{\mathcal E}$ governed by the matrix $\mathbb A(\theta)$, the Floquet matrix and the compact form of that system (lines 524--562); the discriminant corollary that defines $\Delta$ and the growth rate $\sigma$ (lines 756--766); and Theorem 11 with its complete proof for the unstratified case (lines 935--1003). Every cross-reference inside the retained spans resolves inside the delivered document, and the only citation commands of the delivered text are the two of the counted proof, whose entries are transcribed from the paper's own bibliography with the paper's own numbers. Printed numbers are the paper's own: the wrapper restores the paper's global theorem counter and its global equation counter before each retained group, so the retained corollary prints as Corollary 8, the retained unstratified theorem as Theorem 11, the counted statement as Theorem 13, and every numbered display of the retained spans carries the paper's own equation number (15), (28), (40)--(42), (65)--(66), (72)--(79) and (82)--(92). Omitted from this package and not redistributed: the paper's preamble and front matter as such (lines 1--51 and 106--254, of which only the title and the author name are reproduced in the wrapper); the introduction and the physical derivation preceding the short-wavelength analysis; the parts of the analysis outside the retained ranges; the numerical study, the figures and the discussion that follow; and the remaining 48 entries of the paper's own bibliography. The delivered document therefore carries a bibliography of exactly the two works its retained text cites, with the paper's own printed numbers, and no reference or citation command of the delivered text was rewritten or adapted. Printed slips kept literal, among them the paper's own wording and its own choice of symbols. Layout and class: the paper's own class is amsart with the reqno option on a custom 192mm by 262mm paper size; the delivered wrapper is the lane's pinned amsart editorial wrapper at 10pt on A4, to which this package adds the paper's own theorem-like declarations with their shared global counter and its own mathematical shorthand macros, all copied verbatim from the paper's preamble. The paper suppresses the printed numbers of displays that its text does not reference; the wrapper leaves them printed, so unreferenced displays of the retained spans additionally show the same numbers the paper assigns them. That difference is editorial and changes no numbering. No publisher class file, no publisher style file and no upstream script is redistributed. Assets: no retained span of this package contains a figure environment, a drawing picture or a graphics-inclusion command, no image or artwork file exists in or is redistributed with this package, and the manifest and the delivery evidence both carry an empty asset-dependency list; the paper's own figures lie outside every retained span and are not copied. A full rights notice is delivered at licenses/arxiv-2604.03620-CC-BY-4.0.txt.
\setcounter{equation}{14}
\begin{equation}\label{isentropic}
   P=K\,\rho^\gamma, \qquad\text{where} \qquad \gamma=\frac{1}{1-\kappa},
\end{equation}

\setcounter{equation}{27}
  \begin{equation}\label{leading order}
   \left \{\begin{aligned}
    \dot{\x}&=\bu(\x,t),\\
     \dot{\boldsymbol{\xi}}& =-(\nabla\bu)^T\boldsymbol{\xi}, \\  
    \dot{\A}&=-(\A\cdot\nabla)\bu+2\dfrac{\boldsymbol{\xi}\cdot[(\A\cdot\nabla)\bu]}{||\boldsymbol{\xi}||^2}\boldsymbol{\xi}+\cs\left(\nabla g_1-\dfrac{\nabla g_1\cdot\boldsymbol{\xi}}{||\boldsymbol{\xi}||^2}\boldsymbol{\xi}\right)\B, \\
       \dot{\B}&=-\A\cdot\cs \nabla g_3 ,
       \end{aligned} \right.
\end{equation}

\setcounter{equation}{39}
\begin{itemize}
\item $\nu^2>0$: stable (statically stable) stratification;
\item $\nu^2=0$: neutral stratification -- a well-mixed, dry-adiabatic layer, $K=$ const in \eqref{isentropic};
\item $\nu^2<0$: statically unstable stratification.
\end{itemize}
With $\boldsymbol{\mathcal E}=(\mathfrak{U}_1,\mathfrak{U}_3,\widetilde{\B})^T$ the system becomes
\begin{equation}\label{eq:Asys}
\boldsymbol{\mathcal E}'=\mathbb A(\theta)\boldsymbol{\mathcal E},\qquad\text{with}\qquad
\mathbb A(\theta)=
\begin{pmatrix}
-\dfrac{\ee\sin\theta}{\D} & \dfrac{\ee(\cos\theta-\ee)}{\D} & \ee\sin\theta\\[1.2ex]
\dfrac{\ee(\cos\theta+\ee)}{\D} & \dfrac{\ee\sin\theta}{\D} & 1-\ee\cos\theta\\[1.2ex]
-\nu^2\dfrac{\ee\sin\theta}{\D} & -\nu^2\dfrac{1-\ee\cos\theta}{\D} & 0
\end{pmatrix},
\end{equation}
and we have denoted by $'={\rm d}/{\rm d}\theta$.

\begin{remark}\label{rem:params}
The uniform velocities $\bar{u},\bar{w}$ do not appear in \eqref{eq:Asys}, and the Mach number $c/\cs$ enters only through $\nu^2$. For a given level $b$ the stability problem therefore depends on two numbers only: the steepness $\ee=e^{kb}\in(0,1)$ of the trochoidal isoline $b=\mathrm{const}$ (whose amplitude is $e^{kb}/k$), and the stratification parameter $\nu^2$.
\end{remark}
\noindent
It is convenient to write $\mathbf V=(\mathfrak{U}_1,\mathfrak{U}_3)^T$ and to introduce the $2\pi$-periodic vectors
\begin{equation}\label{eq:pq}
\mathbf p(\theta)=\begin{pmatrix}1-\ee\cos\theta\\-\ee\sin\theta\end{pmatrix}=\begin{pmatrix}
    x_a\\ z_a
\end{pmatrix},\qquad
\mathbf q(\theta)=\begin{pmatrix}\ee\sin\theta\\1-\ee\cos\theta\end{pmatrix}=J\mathbf p=\D\,\begin{pmatrix}
    b_x\\ b_z
\end{pmatrix},
\end{equation}
where $\displaystyle J=\begin{pmatrix}0&-1\\1&0\end{pmatrix}$. The vectors $\mathbf p$ and $\mathbf q$ are tangent and normal, respectively, to the isolines $b=\mathrm{const}$. They satisfy $\mathbf p\cdot\mathbf q=0$ and $|\mathbf p|^2=|\mathbf q|^2=h(\theta):=1-2\ee\cos\theta+\ee^2>0$. Then \eqref{eq:Asys} reads
\begin{equation}\label{eq:compact}
\begin{aligned}
    \mathbf V'&=\Msf\,\mathbf V+\widetilde{\B}\,\mathbf q,\\
\widetilde{\B}'&=-\frac{\nu^2}{\D}\,\mathbf q\cdot\mathbf V,
\end{aligned}\qquad
\Msf:=\frac{1}{kc}\begin{pmatrix}u_x&u_z\\w_x&w_z\end{pmatrix}
=\frac{\ee}{\D}\begin{pmatrix}-\sin\theta&\cos\theta-\ee\\ \cos\theta+\ee&\sin\theta\end{pmatrix}.
\end{equation}

\setcounter{theorem}{7}
\setcounter{equation}{64}
\begin{corollary}\label{cor:discriminant}
Let $\ee\in(0,1)$ and let $\{1,\mu,\mu^{-1}\}$ be the Floquet multipliers of \eqref{eq:Asys}. Set
\begin{equation}\label{eq:Delta}
\Delta(\ee,\nu^2):=\tr\mathbb M-1=\mu+\mu^{-1}\in\mathbb R .
\end{equation}
Then the amplitudes $[\mathfrak{U}_1,\mathfrak{U}_3,\widetilde{\B}]$ (equivalently $(\A,\B)$) grow exponentially in $t$ if and only if $|\Delta|>2$, in which case the growth rate is
\begin{equation}\label{eq:sigma}
\sigma=\frac{\Omega}{2\pi}\,\arcosh\frac{|\Delta|}{2}.
\end{equation}
If $|\Delta|<2$ the amplitudes remain bounded, while if $|\Delta|=2$ there is no exponential growth, but the amplitudes may grow algebraically.
\end{corollary}

\setcounter{theorem}{10}
\setcounter{equation}{71}

\subsubsection{Neutral stratification: an explicit solution}\label{sec_neutral}
When $\nu=0$ the amplitude $\widetilde{\B}$ is constant, the monodromy matrix is block triangular, $\mathbb M=\left(\begin{smallmatrix}X(2\pi)&*\\0&1\end{smallmatrix}\right)$ with $X$ the fundamental matrix of $\mathbf V'=\Msf\mathbf V$, and therefore $\Delta=\tr X(2\pi)$: the relevant dynamics is $\mathbf V'=\Msf\mathbf V$. Writing
\[
R_\alpha=\begin{pmatrix}
    \cos\alpha & -\sin\alpha\\
    \sin\alpha &\cos\alpha
\end{pmatrix}
\]
for the rotation by an angle $\alpha$, one checks that
\[
\Msf=\frac{\ee}{\D}\,R_{\alpha}\begin{pmatrix}1&0\\0&-1\end{pmatrix}R_{\alpha}^T+\frac{\ee^2}{\D}\,J,\qquad \alpha(\theta)=\frac{\theta}{2}+\frac{\pi}{4},
\]
that is, the principal axes of strain rotate at exactly half the orbital angular velocity. In the co-rotating frame $\mathbf V=R_{\alpha(\theta)}\mathbf W$ the system becomes autonomous:
\begin{equation}\label{eq:W}
\mathbf W'=\mathsf K\,\mathbf W,\qquad
\mathsf K=\begin{pmatrix}\eta&-\varpi\\ \varpi&-\eta\end{pmatrix},
\end{equation}
where
\begin{equation}\label{eq:gsome}
\eta:=\frac{\ee}{\D},\qquad\text{and}\qquad
\varpi:=\frac{\ee^{2}}{\D}-\frac12=\frac{3\ee^{2}-1}{2\D},
\end{equation}
the latter being the vorticity of the basic flow corrected by the rotation $\frac{1}{2}$ of the frame itself. The eigenvalues of $\mathsf K$ are $\pm\lambda$ with
\begin{equation}\label{eq:lambda}
\lambda^2=\eta^2-\varpi^2=\frac{9\ee^2-1}{4(1-\ee^2)},
\end{equation}
and the fundamental matrix of the system \eqref{eq:W} is $e^{\theta\mathsf K}$. From 
\begin{equation}
    \mathbf V(\theta)=R_{\alpha(\theta)}\mathbf W(\theta)=R_{\alpha(\theta)}e^{\theta\mathsf K}\mathbf W(0) =R_{\alpha(\theta)}e^{\theta\mathsf K}R_{\pi/4}^T \mathbf V(0)
\end{equation}
we note that the fundamental matrix of $\mathbf V'=\Msf\mathbf V$ is $X(\theta)=R_{\theta/2+\pi/4}\,e^{\theta\mathsf K}R_{\pi/4}^T$, so that the monodromy matrix $X(2\pi)=-R_{\pi/4}e^{2\pi\mathsf K}R_{\pi/4}^T$; its trace equals $-\tr e^{2\pi\mathsf K}$ by invariance of the trace under similarity. Finally, $\tr\mathsf K=0$ and $\det\mathsf K=-\lambda^2$, so the Cayley--Hamilton theorem gives $\mathsf K^2=\lambda^2\,\mathbb I$ and hence
\[
e^{\theta\mathsf K}=\cosh(\lambda\theta)\,\mathbb I+\frac{\sinh(\lambda\theta)}{\lambda}\,\mathsf K.
\]
Taking the trace at $\theta=2\pi$ we obtain
\begin{equation}\label{eq:Delta0}
\Delta(\ee,0)=-2\cosh(2\pi\lambda).
\end{equation}
The nature of the eigenvalues \eqref{eq:lambda} determines the behaviour of $\mathbf{W}$ (hence, of $\mathbf{V}$). If $\ee>1/3$, \eqref{eq:lambda} has two real solution---one positive and one negative---giving an exponential growth of the perturbation, leading to the instability of the problem. This is summarised in the following Theorem.

\begin{theorem}\label{thm:neutral}
Let $\nu=0$ and $\ee=e^{kb}>1/3$, so that the eigenvalues of $\mathsf K$ are 
\begin{equation}
    \pm\lambda\in \mathbb{R} \qquad\text{with}\qquad\lambda=\frac{1}{2} \sqrt{\frac{9\ee^2-1}{1-\ee^2}}>0.
\end{equation}
We have the following
\begin{enumerate}
    \item[(i)] the short-wavelength perturbation at level $b$ grows exponentially, with growth rate
\begin{equation}\label{eq:sigma0}
\sigma=\Omega\,\lambda=\frac{\Omega}{2}\sqrt{\frac{9\ee^2-1}{1-\ee^2}},\qquad \Omega=\sqrt{gk};\end{equation}
\item[(ii)] $\Delta<-2$, namely the growing mode is subharmonic: it changes sign every orbital period, oscillating at $\Omega/2$;
\item[(iii)] the choice $\mathfrak{U}_2\equiv0$, $\B\equiv0$,
\begin{equation}\label{eq:explicit}
\begin{pmatrix}\mathfrak{U}_1\\\mathfrak{U}_3\end{pmatrix}(t)
=e^{-\lambda\,\theta(t)}\,R_{\theta(t)/2+\pi/4}\begin{pmatrix}\varpi\\\eta+\lambda\end{pmatrix},
\qquad \theta(t)=ka-\Omega t,
\end{equation}
is an exact solution of the WKB amplitude system \eqref{leading order} (restricted to $\boldsymbol\xi_0=(0,1,0)$) along the particle path with label $b=k^{-1}\ln\ee$, and $|(\mathfrak{U}_1,\mathfrak{U}_3)(t)|=|(\varpi,\eta+\lambda)|\,e^{\lambda(\Omega t-ka)}$ grows without bound as $t\to\infty$. 
\end{enumerate}
\end{theorem}
\begin{proof}
\begin{enumerate}
    \item[(i)] follows from \eqref{eq:Delta0} and Corollary~\ref{cor:discriminant}: for $\ee>1/3$, $\lambda>0$ is real and $|\Delta|=2\cosh(2\pi\lambda)>2$, giving \eqref{eq:sigma0} by \eqref{eq:sigma}.
\item[(ii)] the multipliers $-e^{\pm2\pi\lambda}$ are negative whenever $\ee>1/3$, giving the subharmonic character as in Corollary~\ref{cor:discriminant}.
\item[(iii)] An easy direct computation shows that the vector $\mathbf{z}=(\varpi,\ \eta+\lambda)^T$ satisfies $\mathsf K\mathbf{z}=-\lambda\mathbf{z}$.\\
By construction $\mathbf W(\theta)=e^{-\lambda\theta}(\varpi,\eta+\lambda)^T$ solves $\mathbf W'=\mathsf K\mathbf W$, being a multiple of the eigenvector for the eigenvalue $-\lambda$; then $\mathbf V(\theta)=R_{\alpha(\theta)}\mathbf W(\theta)$ solves $\mathbf V'=\Msf\mathbf V$ by the derivation of \eqref{eq:W}. Since $\theta$ decreases linearly with $t$ at rate $\Omega$, the factor $e^{-\lambda\theta(t)}$ grows like $e^{\lambda\Omega t}$, and $R_\alpha$, being a rotation, does not affect the modulus.
\end{enumerate}
\end{proof}

\setcounter{theorem}{12}
\setcounter{equation}{81}
\begin{theorem}\label{prop:strat}
$\Delta(\ee,\nu^2)$, defined by \eqref{eq:Delta}, is real-analytic on $(0,1)\times\mathbb R$. Moreover:
\begin{enumerate}
\item[(i)] the unstable set $\mathcal S:=\{|\Delta|>2\}$ is open and contains the segment $\{(\ee,0):\ee>1/3\}$; hence for every $\ee\in(1/3,1)$ there is $\nu_*(\ee)>0$, possibly small, such that the level is unstable for all $|\nu^2|<\nu_*(\ee)^2$;
\item[(ii)] there are a neighbourhood of $(\ee,\nu^2)=(1/3,0)$ and a real-analytic function $\ee_c$ with $\ee_c(0)=1/3$ such that, in that neighbourhood, the level is unstable if and only if $\ee>\ee_c(\nu^2)$, and
\begin{equation}\label{eq:ec}
\ee_c(\nu^2)=\frac{1}{3}-\frac{4}{9}\,\nu^2+\mathcal{O}(\nu^4);
\end{equation}
\item[(iii)] at $\ee=0$ one has $\Delta=2\cos(2\pi\nu)$ for $\nu^2\ge0$ and $\Delta=2\cosh(2\pi|\nu|)$ for $\nu^2<0$; hence statically unstable layers (i.e. $\nu^2<0$) are unstable already for arbitrarily small steepness, while for $\nu^2>0$ the closure of $\mathcal S$ meets the axis $\ee=0$ only at the resonant values $\nu=n/2$, $n\in\mathbb N$.
\end{enumerate}
\end{theorem}

\begin{proof}
\begin{enumerate}
    \item [(i)] Analyticity of $\Delta$ follows from the analytic dependence of $\mathbb A(\theta;\ee,\nu^2)$ on the parameters and from the standard theory of linear ODEs with parameters (the fundamental matrix, hence its trace at $\theta=2\pi$, depends analytically on the coefficients); note that $\mathbb A$ is regular at $\ee=0$ as well. Since $\{|x|>2\}$ is open in $\mathbb R$ and $\Delta$ is continuous, $\mathcal S=\Delta^{-1}(\{|x|>2\})$ is open; it contains $\{\ee>1/3,\ \nu=0\}$ by Theorem~\ref{thm:neutral}. Openness gives, around each such point, a small disc contained in $\mathcal S$, in particular a vertical segment in $\nu^2$.
\item[(ii)] The neutral curve (i.e. the boundary between stable and unstable levels) near $(1/3,0)$ is the set where $\Delta=-2$. Since $\Delta(1/3,0)=-2$ by \eqref{eq:Delta0}, the implicit function theorem applied to $\Delta(\ee,\nu^2)+2=0$ describes this set as a graph $\ee=\ee_c(\nu^2)$, provided $\partial_\ee\Delta(1/3,0)\neq0$, and gives $\ee_c'(0)=-\partial_{\nu^2}\Delta/\partial_\ee\Delta$ at $(1/3,0)$. We compute the two derivatives in closed form.
\begin{itemize}
\item \emph{Derivative in $\ee$.} By \eqref{eq:lambda}--\eqref{eq:Delta0}, $\Delta(\ee,0)=F\big(\lambda^2(\ee)\big)$ for all $\ee\in(0,1)$, where
\[
F(x):=-2\cosh\big(2\pi\sqrt x\big)=-2\sum_{n\ge0}\frac{(2\pi)^{2n}}{(2n)!}\,x^{n},\qquad \lambda^2(\ee)=\frac{9\ee^2-1}{4\D}.
\]
Since $F$ is an entire function of $x$, this representation holds on both sides of $\ee=1/3$, where $\lambda^2$ changes sign, and no branch point occurs. We have $F'(0)=-4\pi^2$, and since the numerator of $\lambda^2$ vanishes at $\ee=1/3$, where $\D=8/9$,
\[
\frac{{\rm d}\lambda^2}{{\rm d}\ee}\Big|_{\ee=\frac{1}{3}}=\frac{18\ee}{4\D}\Big|_{\ee=\frac{1}{3}}=\frac{27}{16}.
\]
Hence
\begin{equation}\label{eq:dDde}
\partial_\ee\Delta(1/3,0)=-4\pi^2\cdot\frac{27}{16}=-\frac{27\pi^2}{4}\neq0 .
\end{equation}




\item \emph{Derivative in $\nu^2$.} The fundamental matrix $\F$ solves $\F'=\mathbb A\F$, $\F(0)=\mathbb I$, and depends on $\nu^2$. Differentiating this problem with respect to $\nu^2$ by the product rule, $\Psi:=\partial_{\nu^2}\F$ solves
\[
\Psi'=\mathbb A\,\Psi+\big(\partial_{\nu^2}\mathbb A\big)\F,\qquad \Psi(0)=0,
\]
the initial value being zero because $\F(0)=\mathbb I$ does not depend on $\nu^2$. This is the equation for $\F$ with the  forcing $\boldsymbol{\mathfrak{ F}}:=(\partial_{\nu^2}\mathbb A)\F$, and it is solved by the variation-of-constants formula: the solution of $\mathbf Y'=\mathbb A\mathbf Y+\boldsymbol{\mathfrak{ F}}(\theta)$ with $\mathbf Y(0)=0$ is
\[
\mathbf Y(\theta)=\F(\theta)\int_0^{\theta}\F(s)^{-1}\boldsymbol{\mathfrak{ F}}(s)\,{\rm d}s,
\]
as one checks by direct computations. Evaluating at $\theta=2\pi$ and $\nu=0$, and writing $\F_0:=\F|_{\nu=0}$, we obtain
\[
\partial_{\nu^2}\mathbb M\big|_{\nu=0}=\F_0(2\pi)\int_0^{2\pi}\F_0(s)^{-1}\,\partial_{\nu^2}\mathbb A(s)\,\F_0(s)\,{\rm d}s .
\]
Since $\Delta=\tr\mathbb M-1$ and the trace is linear, $\partial_{\nu^2}\Delta=\tr(\partial_{\nu^2}\mathbb M)$, and the trace may be taken inside the integral. By \eqref{eq:Asys}, only the entries $(3,1)$ and $(3,2)$ of $\mathbb A$ depend on $\nu^2$, and linearly; hence
\[
\partial_{\nu^2}\mathbb A(s)=-\frac{1}{\D}\begin{pmatrix}0&0&0\\0&0&0\\ q_1(s)&q_2(s)&0\end{pmatrix}
=-\frac{1}{\D}\,\mathbf e_3\,\widetilde{\mathbf q}(s)^T,
\]
where $\mathbf e_3=(0,0,1)^T$ and $\widetilde{\mathbf q}:=(q_1,q_2,0)^T$: a matrix of rank one. For a column vector $\mathbf a$ and a row vector $\mathbf b^T$, the cyclicity of the trace gives $\tr(\mathbf a\,\mathbf b^T)=\mathbf b^T\mathbf a$. Applying this with $\mathbf a=\F_0(2\pi)\F_0(s)^{-1}\mathbf e_3$ and $\mathbf b^T=\widetilde{\mathbf q}(s)^T\F_0(s)$ yields
\begin{equation}\label{eq:dDdnu0}
\partial_{\nu^2}\Delta\big|_{\nu=0}=-\frac{1}{\D}\int_0^{2\pi}\widetilde{\mathbf q}(s)^T\,\F_0(s)\,\F_0(2\pi)\,\F_0(s)^{-1}\,\mathbf e_3\,{\rm d}s .
\end{equation}

Let us describe now the structure of $\F_0$. At $\nu=0$ the last equation in \eqref{eq:compact} reads $\widetilde{\B}'=0$, so $\widetilde{\B}$ is constant and enters the first equation, $\mathbf V'=\Msf\mathbf V+\widetilde{\B}\,\mathbf q$, as a constant forcing coefficient. Denoting by $X(\theta)$ the fundamental matrix of $\mathbf V'=\Msf\mathbf V$, the variation-of-constants formula gives
\[
\mathbf V(\theta)=X(\theta)\Big[\mathbf V(0)+\widetilde{\B}(0)\,\mathbf j(\theta)\Big],\qquad
\mathbf j(\theta):=\int_0^{\theta}X(s)^{-1}\mathbf q(s)\,{\rm d}s,
\]
that is, in the variables $(\mathbf V,\widetilde{\B})$,
\[
\F_0(\theta)=\begin{pmatrix}X(\theta)&X(\theta)\,\mathbf j(\theta)\\ 0&1\end{pmatrix},\qquad
\F_0(\theta)^{-1}=\begin{pmatrix}X(\theta)^{-1}&-\mathbf j(\theta)\\ 0&1\end{pmatrix},
\]
the inverse being checked by multiplication, since the upper right block of the product is $-X\mathbf j+X\mathbf j=0$. Hence $\F_0(s)^{-1}\mathbf e_3=(-\mathbf j(s),1)^T$, and successively
\[
\begin{aligned}
\F_0(2\pi)\begin{pmatrix}-\mathbf j(s)\\ 1\end{pmatrix}&=\begin{pmatrix}X(2\pi)\big(\mathbf j(2\pi)-\mathbf j(s)\big)\\ 1\end{pmatrix},\\[4pt]
\F_0(s)\begin{pmatrix}X(2\pi)\big(\mathbf j(2\pi)-\mathbf j(s)\big)\\ 1\end{pmatrix}
&=\begin{pmatrix}X(s)\big[X(2\pi)\big(\mathbf j(2\pi)-\mathbf j(s)\big)+\mathbf j(s)\big]\\ 1\end{pmatrix}.
\end{aligned}
\]
Since the third component of $\widetilde{\mathbf q}$ vanishes, only the first two components contribute in \eqref{eq:dDdnu0}, and
\begin{equation}\label{eq:dDdnu}
\partial_{\nu^2}\Delta\big|_{\nu=0}=-\frac1\D\int_0^{2\pi}\mathbf q(s)\cdot X(s)\Big[X(2\pi)\big(\mathbf j(2\pi)-\mathbf j(s)\big)+\mathbf j(s)\Big]\,{\rm d}s .
\end{equation}

By the construction leading to \eqref{eq:W}--\eqref{eq:Delta0}, the change of variables $\mathbf V=R_{\alpha(\theta)}\mathbf W$, with $\alpha(\theta)=\theta/2+\pi/4$, transforms $\mathbf V'=\Msf\mathbf V$ into $\mathbf W'=\mathsf K\mathbf W$, so that $\mathbf W(\theta)=E(\theta)\mathbf W(0)$ with $E(\theta):=e^{\theta\mathsf K}$. Since $\alpha(0)=\pi/4$ and the inverse of a rotation is its transpose, $\mathbf W(0)=R_{\pi/4}^T\mathbf V(0)$, and therefore
\[
X(\theta)=R_{\alpha(\theta)}E(\theta)R_{\pi/4}^T,\qquad X(\theta)^{-1}=R_{\pi/4}E(\theta)^{-1}R_{\alpha(\theta)}^T .
\]
At $\theta=2\pi$ one has $\alpha=\pi+\pi/4$, and a further rotation by $\pi$ changes the sign, $R_{\pi+\beta}=-R_\beta$; hence $X(2\pi)=-R_{\pi/4}E(2\pi)R_{\pi/4}^T$. Set
\begin{equation}\label{J}
\mathbf r(s):=R_{\alpha(s)}^T\mathbf q(s),\qquad \boldsymbol{\mathfrak{J}}(s):=\int_0^{s}E(\sigma)^{-1}\mathbf r(\sigma)\,{\rm d}\sigma,
\end{equation}
\begin{equation}\label{ys}
\mathbf y(s):=\boldsymbol{\mathfrak{J}}(s)-E(2\pi)\big(\boldsymbol{\mathfrak{J}}(2\pi)-\boldsymbol{\mathfrak{J}}(s)\big).
\end{equation}
From the expression of $X^{-1}$ we have $\mathbf j=R_{\pi/4}\boldsymbol{\mathfrak{J}}$, so that the bracket in \eqref{eq:dDdnu} equals $R_{\pi/4}\mathbf y(s)$, the minus sign in front of $E(2\pi)$ coming from $X(2\pi)$. Then $X(s)R_{\pi/4}\mathbf y=R_{\alpha(s)}E(s)\mathbf y$, and since $\mathbf q\cdot(R\mathbf v)=(R^T\mathbf q)\cdot\mathbf v$ for every rotation $R$, \eqref{eq:dDdnu} becomes
\begin{equation}\label{eq:dDdnurot}
\partial_{\nu^2}\Delta\big|_{\nu=0}=-\frac1\D\int_0^{2\pi}\mathbf r(s)\cdot E(s)\,\mathbf y(s)\,{\rm d}s .
\end{equation}

It is now convenient to introduce the following orthonormal vectors $\boldsymbol{\upsilon}_\pm:=(1,\pm1)^T/\sqrt2$; if vectors are written in coordinates, $\mathbf a=a_+\boldsymbol{\upsilon}_++a_-\boldsymbol{\upsilon}_-$, the scalar product is $\mathbf a\cdot\mathbf b=a_+b_++a_-b_-$. With $\mathsf K=\left(\begin{smallmatrix}\eta&-\varpi\\ \varpi&-\eta\end{smallmatrix}\right)$ from \eqref{eq:W}, a direct multiplication gives
\[
\mathsf K\boldsymbol{\upsilon}_+=(\eta-\varpi)\,\boldsymbol{\upsilon}_-,\qquad \mathsf K\boldsymbol{\upsilon}_-=(\eta+\varpi)\,\boldsymbol{\upsilon}_+ .
\]
To express $\mathbf r$, put $u=s/2$, so that $\alpha=u+\pi/4$, and $(r_1,r_2)=R_\alpha^T\mathbf q=(q_1\cos\alpha+q_2\sin\alpha,\ -q_1\sin\alpha+q_2\cos\alpha)$ with $\mathbf q=(\ee\sin s,\,1-\ee\cos s)$. Using $\cos\alpha-\sin\alpha=-\sqrt2\sin u$ and $\cos\alpha+\sin\alpha=\sqrt2\cos u$, together with $\cos(s-u)=\cos u$ and $\sin(s-u)=\sin u$ (because $s-u=u$), we find
\[
\begin{aligned}
r_1+r_2&=\sqrt2\,\big(q_2\cos u-q_1\sin u\big)=\sqrt2\,\big(\cos u-\ee\cos(s-u)\big)=\sqrt2\,(1-\ee)\cos u,\\
r_1-r_2&=\sqrt2\,\big(q_1\cos u+q_2\sin u\big)=\sqrt2\,\big(\sin u+\ee\sin(s-u)\big)=\sqrt2\,(1+\ee)\sin u,
\end{aligned}
\]
that is, since $r_\pm=(r_1\pm r_2)/\sqrt2$,
\begin{equation}
    \label{r}
\mathbf r(s)=\underbrace{(1-\ee)\cos\frac{s}{2}}_{:=r_+}\,\boldsymbol{\upsilon}_++\underbrace{(1+\ee)\sin\frac{s}{2}}_{:=r_-}\,\boldsymbol{\upsilon}_- .
\end{equation}

Now set $\ee=1/3$, giving $\D=8/9$, $\eta=\ee/\D=3/8$ and $\varpi=(3\ee^2-1)/(2\D)=-3/8$ by \eqref{eq:gsome}, so that
\[
\mathsf K\boldsymbol{\upsilon}_+=\frac{3}{4}\,\boldsymbol{\upsilon}_-,\qquad \mathsf K\boldsymbol{\upsilon}_-=0 .
\]
Applying $\mathsf K$ twice to each basis vector gives zero, i.e.\ $\mathsf K^2=0$, and the exponential series terminates: $E(s)=\mathbb I+s\mathsf K$. In coordinates,
\begin{equation}
    \begin{aligned}
        &E(s)(x_+,x_-)=\big(x_+,\ x_-+\frac{3}{4} s\,x_+\big),\\
&E(s)^{-1}(x_+,x_-)=\big(x_+,\ x_--\frac{3}{4} s\,x_+\big),
    \end{aligned}
\end{equation}
and in particular $E(2\pi)(x_+,x_-)=\big(x_+,\ x_-+\frac{3\pi}{2}x_+\big)$. \\
Using \eqref{r}, it follows that
\[
\mathbf r(s)=\Big(\frac{2}{3}\cos\frac{s}{2},\ \frac{4}{3}\sin\frac{s}{2}\Big),\qquad
E(s)^{-1}\mathbf r(s)=\Big(\frac{2}{3}\cos\frac{s}{2},\ \frac{4}{3}\sin\frac{s}{2}-\frac{s}{2}\,\cos\frac{s}{2}\Big).
\]

Write $\boldsymbol{\mathfrak{J}}=(\mathfrak{j}_+, \mathfrak{j}_-)$. Computing the first and second components gives:
\[
\begin{aligned}
\mathfrak{j}_+(s)&=\int_0^s\frac{2}{3}\cos\frac{\sigma}{2}\,{\rm d}\sigma=\frac{4}{3}\sin\frac{s}{2},\\
\mathfrak{j}_-(s)
&=\int_0^s\Big( \frac{4}{3}\sin\frac{\sigma}{2}-\frac{1}{2}\,\sigma\cos\frac{\sigma}{2}\Big){\rm d}\sigma=\frac{14}{3}\Big(1-\cos\frac{s}{2}\Big)-s\sin\frac{s}{2} .
\end{aligned}
\]
In particular $\boldsymbol{\mathfrak{J}}(2\pi)=\big(0,\frac{28}{3}\big)$, and
\[
\boldsymbol{\mathfrak{J}}(2\pi)-\boldsymbol{\mathfrak{J}}(s)=\Big(-\frac{4}{3}\sin\frac{s}{2},\ \ \frac{14}{3}+\frac{14}{3}\cos\frac{s}{2}+s\sin\frac{s}{2}\Big).
\]
Applying $E(2\pi)$, which adds $\frac{3\pi}{2}$ times the first component to the second, and subtracting the result from $\boldsymbol{\mathfrak{J}}(s)$ according to \eqref{ys}, we obtain
\[
\mathbf y(s)=\Big(\frac{8}{3}\sin\frac{s}{2},\ \ -\frac{28}{3}\cos\frac{s}{2}-2s\sin\frac{s}{2}+2\pi\sin\frac{s}{2}\Big),
\]
and applying $E(s)$, which adds $\frac{3}{4} s\,y_+=2s\sin\frac{s}{2}$ to the second component,
\[
E(s)\,\mathbf y(s)=\Big(\frac{8}{3}\sin\frac{s}{2},\ \ 2\pi\sin\frac{s}{2}-\frac{28}{3}\cos\frac{s}{2}\Big).
\]
 Hence, using the duplication formulas $\sin\frac{s}{2}\cos\frac{s}{2}=\frac{1}{2}\sin s$, $\sin^2\frac{s}{2}=\frac{1}{2}(1-\cos s)$, we have
\[
\mathbf r(s)\cdot E(s)\,\mathbf y(s)=-\frac{16}{3}\sin s+\frac{4\pi}{3}\,(1-\cos s).
\]
The integrals of $\sin s$ and $\cos s$ over $[0,2\pi]$ vanish, so the integral in \eqref{eq:dDdnurot} equals $\frac{4\pi}{3}\cdot2\pi=\frac{8\pi^2}{3}$, and \eqref{eq:dDdnurot} with $\D=8/9$ gives
\begin{equation}\label{eq:dDdnu2}
\partial_{\nu^2}\Delta(1/3,0)=-\frac98\cdot\frac{8\pi^2}{3}=-3\pi^2 .
\end{equation}
In conclusion, since $\Delta(1/3,0)+2=0$ and $\partial_\ee\Delta(1/3,0)\neq0$, the implicit function theorem applies, and
\begin{equation}
\ee_c'(0)=-\frac{-3\pi^2}{-\frac{27\pi^2}{4}}=-\frac49.
\end{equation}
As $\ee_c$ is real-analytic in $\nu^2$, its Taylor expansion reads $\ee_c(\nu^2)=\frac13-\frac49\nu^2+\mathcal{O}\big((\nu^2)^2\big)$, which is \eqref{eq:ec}.\\
Finally, near $(1/3,0)$ we have $\Delta$ close to $-2$, hence $\Delta<2$, and $\partial_\ee\Delta<0$ by continuity of the derivative, so that $\Delta$ is strictly decreasing in $\ee$ at fixed $\nu^2$. Consequently $\Delta>-2$, i.e.\ $|\Delta|<2$, for $\ee<\ee_c(\nu^2)$, and $\Delta<-2$, i.e.\ $|\Delta|>2$, for $\ee>\ee_c(\nu^2)$: in that neighbourhood the level is unstable if and only if $\ee>\ee_c(\nu^2)$.
\end{itemize}
\item[(iii)]
Since $\mathbb A$ is regular at $\ee=0$, $\Delta$ is analytic, and in particular continuous, up to $\ee=0$. At $\ee=0$ one has $\Msf\equiv0$ and $\mathbf q\equiv(0,1)^T$, so that \eqref{eq:Asys} has the constant matrix
\[
\mathbb A_0=\begin{pmatrix}0&0&0\\0&0&1\\0&-\nu^2&0\end{pmatrix},
\qquad\text{i.e.}\qquad
\mathfrak{U}_1'=0,\quad \mathfrak{U}_3''=-\nu^2\,\mathfrak{U}_3 .
\]
Physically, $\mathfrak{U}_3$ performs a buoyancy oscillation of frequency $\nu$ in $\theta$, that is of frequency $\nu\Omega=\mathcal N$ in time, while $\mathfrak{U}_1$ is frozen. Hence $\mathbb M=e^{2\pi\mathbb A_0}$ is block diagonal, with the entry $1$ for $\mathfrak{U}_1$ and the block
\[
\begin{cases}
\begin{pmatrix}
\cos (2\pi\nu) & \nu^{-1}\sin (2\pi\nu)\\
-\nu\sin (2\pi\nu) & \cos (2\pi\nu)
\end{pmatrix},
& \quad\text{if}\quad\nu^2>0,  \\[1.2em]
\begin{pmatrix}
\cosh (2\pi|\nu|) & |\nu|^{-1}\sinh (2\pi|\nu|)\\
|\nu|\sinh (2\pi|\nu|) & \cosh (2\pi|\nu|)
\end{pmatrix},
& \quad\text{if}\quad \nu^2<0.
\end{cases}
\]
 Therefore $\tr\mathbb M=1+2\cos(2\pi\nu)$ or $\tr\mathbb M=1+2\cosh(2\pi|\nu|)$ , and
\[
\Delta(0,\nu^2)=2\cos(2\pi\nu)\quad(\nu^2\ge0),\qquad \Delta(0,\nu^2)=2\cosh(2\pi|\nu|)\quad(\nu^2<0),
\]
the two expressions agreeing, with value $2$, at $\nu=0$. Two consequences follow by continuity of $\Delta$.

If $\nu^2<0$, then $\Delta(0,\nu^2)>2$, hence $\Delta(\ee,\nu^2)>2$ for all sufficiently small $\ee\geq0$: a statically unstable layer is unstable for arbitrarily small steepness, with growth rate $\sigma\to\frac{\Omega}{2\pi}\arcosh\big(\cosh2\pi|\nu|\big)=\Omega|\nu|$ as $\ee\to0$, the classical convective growth rate (see e.g. \cite{Chandrasekhar,Vallis}).

If $\nu^2>0$ and $2\nu\notin\mathbb Z$, then $|\Delta(0,\nu^2)|=|2\cos2\pi\nu|<2$ strictly, hence $|\Delta|<2$ on a whole neighbourhood of $(0,\nu^2)$; this neighbourhood does not meet $\mathcal S$, and $(0,\nu^2)$ does not belong to the closure of $\mathcal S$. The only points of the axis $\ee=0$ with $\nu^2>0$ that may belong to that closure are therefore those with $|2\cos2\pi\nu|=2$, i.e.\ $\nu=n/2$ with $n\in\mathbb N$.
\end{enumerate}
\end{proof}

\begin{thebibliography}{99}
\bibitem[5]{Chandrasekhar} S.~Chandrasekhar, \textit{Hydrodynamic and Hydromagnetic Stability}, Dover Publications, New York, 1981.
\bibitem[48]{Vallis} G.~K. Vallis, \textit{Atmospheric and Oceanic Fluid Dynamics: Fundamentals and Large-Scale Circulation}, 2nd ed., Cambridge University Press, Cambridge, 2017.
\end{thebibliography}
\end{document}
